\section{Packing the exterior tensor}\label{sec:packing}

Theorem~\ref{ext:star} chooses a Vandermonde matrix after the row sites
and their prices are known. The circuit lower bound of
Theorem~\ref{thm:lowdegree} requires a single tensor with zero-one
coefficients, specified from the integers $d$ and $w$ alone. We first
group exterior copies into row letters and compress the row alphabet.
The resulting tensor retains the star price to within a constant per
row. We then encode its Vandermonde minors by integers so that every
choice of matrix can be recovered from one binary tensor by linear maps
on the rows.

\subsection{Weighted prices on ordinary tensor powers}
Let $A$ have ordered bases on its $d$ legs, placed at
$x=(x_1,\ldots,x_d)$. Give each basis letter $v$ on leg $i$ a real
weight $t_i(v)$. On that leg of the grouped power $A^{\otimes\ell}$,
a word has weight equal to the sum of its $\ell$ letter weights.
For a weight vector $a=(a_1,\ldots,a_d)$ let $[A^{\otimes\ell}]_a$
be the sector obtained by projecting each leg onto the words of
weight $a_i$. We will use the weighted price
\begin{equation}\label{pack:weighted}
 \mathcal B_x(A)=\lim_{\ell\to\infty}\frac{1}{\ell}
 \max_a\left( H([A^{\otimes\ell}]_a;x)+\sum_{i=1}^d a_i\right).
\end{equation}
The zero sector has price $-\infty$. For $A\ne0$, the unnormalized
maxima are superadditive. Indeed, projecting a sector of combined
weight onto prescribed weights in two consecutive blocks supplies the
product of their sectors. The price of a sector is at most
$\ell H(A;x)$, and its weight is bounded above linearly in $\ell$.
Thus the limit equals the supremum of the normalized maxima. This
plays the role of
the exterior-sector price of Equation~\eqref{ext:value} for ordinary
tensor powers.

\begin{lemma}[Weighted price rules]\label{pack:value}
At any fixed placement, the weighted price has Zariski-closed
sublevels. It is monotone under legwise maps preserving weight spaces
and additive on grouped products when weights on pairs of letters are
added. For tensors $A_1,\ldots,A_L$ in the same weighted format and
$D_x=\delta(\{x_1,\ldots,x_d\})$,
\begin{equation}\label{pack:sum}
 \mathcal B_x\!\left(\sum_{j=1}^L A_j\right)
 \le \max_j\mathcal B_x(A_j)+D_x\log L .
\end{equation}
If $H(A;y)\le Q$ and $r_i$ is the distance from $y_i$ to $x_i$, then
\begin{equation}\label{pack:transport}
 \mathcal B_x(A)\le Q+\sum_{i=1}^d r_i
                 \log\sum_{v\in V_i}\exp(t_i(v)/r_i),
\end{equation}
where $V_i$ is any nonempty alphabet supporting $A$ on leg $i$.
At $r_i=0$ the summand is $\max_{v\in V_i}t_i(v)$.
\end{lemma}

\begin{proof}
For fixed $\ell$ and $a$, the entries of the sector are polynomial in
the entries of $A$. The supremum formula expresses a weighted price
sublevel as the intersection of inverse images of ordinary price
sublevels, which are closed by Lemma~\ref{price:rules}. A
weight-preserving map acts separately on each sector, giving
monotonicity.

In a power of a grouped product $A\otimes B$, refine a sector by
specifying how much of each weight comes from $A$ and how much comes
from $B$. Each refined sector is a product of the two corresponding
sectors. It is a restriction of the original sector. Conversely, the
original sector is the sum of its refinements embedded in the common
format. The number of refinements grows polynomially in $\ell$ for
fixed alphabets, since the letter counts determine the weights. By
Equation~\eqref{price:sum}, the price of the original sector therefore
exceeds the largest refinement price by at most $O(\log(\ell+1))$.
After division by $\ell$ this gives the upper bound for additivity.
Choosing a maximizing sector for each factor gives the lower bound.

For Equation~\eqref{pack:sum}, expand the $\ell$th power of the sum.
In each of the $L^\ell$ terms, reorder the factors to group identical
ones, and refine a given sector by prescribing the weights within each
group. There are polynomially many refinements per term for fixed $L$
and fixed alphabets. If $A_j$ occurs $\ell_j$ times, the refinement
price plus its total weight is at most
$\sum_{j:\ell_j>0}\ell_j\mathcal B_x(A_j)$ by the supremum formula.
The sum over terms and refinements costs at most
$D_x\ell\log L+O(\log(\ell+1))$ in addition, by
Equation~\eqref{price:sum}. Divide by $\ell$ and take the limit.

For transport, restrict $A^{\otimes\ell}$ at $y$ to a sector of
weight $a$. For $r_i>0$ the number of words of weight $a_i$ on its
$i$th leg is at most
\[
 e^{-a_i/r_i}\left(\sum_{v\in V_i}e^{t_i(v)/r_i}\right)^\ell .
\]
A wire carrying those words from $y_i$ to $x_i$ costs at most $r_i$
times the logarithm of this number. Adding back $a_i$ cancels the
factor $e^{-a_i/r_i}$. For $r_i=0$ the wire is free and
$a_i\le\ell\max_{v\in V_i}t_i(v)$. Taking the maximum and dividing by
$\ell$ proves Equation~\eqref{pack:transport}.
\end{proof}

\subsection{A binary tensor for every placement}

\begin{theorem}[Packed binary witness]\label{pack:witness}
For every sufficiently large even $d$, let $m=\lceil\log_2 d\rceil$
and let $w$ be an integer with $dm\le w\le d^2$. There are a $d$-leg
tensor $T$ with entries in $\{0,1\}$, an alphabet $V$ with
nonpositive weights $t(v)$, and a function $\mathcal F$ on
$(0,\infty)$, all depending only on $d,w$. The tensor has common side
$q_0$ with $w\le\log q_0\le Cw$. Its entries are computable in time
polynomial in $w$.

Let $x_1,\ldots,x_d$ be unit sites with uniform Cartan barycenter at
the tip of a Hadamard cone carrying homogeneous spectral prices over an
uncountable algebraically closed field $K$. Assume the sites have
spare Euclidean coordinates as in Lemma~\ref{price:diagonal}. Over
$K$, the tensor $T$ admits a rowwise restriction to a tensor
$\widehat U$ on the alphabet $V$ satisfying
\begin{align}
 d\mathcal F(1)-Cd &\le \mathcal B_x(\widehat U)\le d\mathcal F(1),\label{pack:high}\\
 r\log\sum_{v\in V} e^{t(v)/r}&\le \mathcal F(r) \qquad(r>0),\label{pack:partition}\\
 \mathcal F(1)&\le Cw,\quad c w\le \mathcal F'(1)\le Cw,\quad
 \sup_{1/2\le r\le2}|\mathcal F''(r)|\le Cw .\label{pack:derivatives}
\end{align}
The constants $c,C>0$ are absolute. The restriction may depend on the
placement and its prices.
\end{theorem}

In the circuit comparison, $\mathcal F(r)$ bounds the cost of
transporting one weighted row a distance $r$. Thus
Equation~\eqref{pack:high} supplies the lower bound at unit distance,
while Equation~\eqref{pack:derivatives} controls the saving when the
two ends of a cut move away from the tip.

\begin{proof}
\emph{Row letters and their weights.}
Take $q$ slots at each row, group $b$ exterior copies into one batch,
and choose the bias range $[-M,M]$ by
\[
 q=\lceil w/m^2\rceil,\qquad n=dq,\qquad b=m^3,\qquad M=500m .
\]
The total number $n$ of slots is even and $M\ge100\log n$. Give slot
$j$ of each row bias $\theta_j=-M+(2j-1)M/q$. By
Theorem~\ref{ext:star} we can choose a Vandermonde matrix $C_V$ for
these sites and biases. Group the slots of $U=\Omega_b(C_V)$ by row.
This leaves the price unchanged, since the slots at a row occupy one
site. A row letter $v\in \{0,1\}^{q\times b}$ records whether slot
$j$ is occupied in copy $e$ by the bit $v_{je}$. Give it weight
\begin{equation}\label{pack:letters}
 t(v)=\sum_{j=1}^q\sum_{e=1}^b
           \big(\theta_j v_{je}-\max(\theta_j,0)\big).
\end{equation}
At a nonzero bias, the favored bit has weight zero and the opposite
bit costs $|\theta_j|$. At zero bias either bit is free. For the full row alphabet the
transport term in Equation~\eqref{pack:transport} factors over the bits,
giving
\begin{equation}\label{pack:free}
 \mathcal F(r)=r\log\sum_{v\in\{0,1\}^{q\times b}}e^{t(v)/r}
      =b\sum_{j=1}^q r\log(1+e^{-|\theta_j|/r}).
\end{equation}

The choice of scales makes $\mathcal F'(1)$ of order $w$. Indeed, for
$z\ge0$ and $1/2\le r\le2$, the summand $r\log(1+e^{-z/r})$ and
its first two derivatives in $r$ are bounded in absolute value by
$C e^{-z/4}$. Its first derivative at one is
$\log(1+e^{-z})+z/(1+e^z)$, bounded below by an absolute positive
constant for $z\le1$. The bias grid has spacing $2M/q$, so
$\sum_j e^{-|\theta_j|/4}=O(1+q/M)$ and $\Omega(q/M)$ slots have
$|\theta_j|\le1$. Since $q/M\to\infty$ and $bq/M=\Theta(w)$, this
proves Equation~\eqref{pack:derivatives} and gives $\mathcal F(2)=O(w)$.

\smallskip
\noindent\emph{From exterior sectors to ordinary powers.}
Compare $\Omega_{b\ell}(C_V)$ with $U^{\otimes\ell}$, using $\ell$
batches of $b$ copies each and ordering the slots by row. Both
coefficients vanish unless each copy has $n/2$ occupied slots.
Otherwise both are a product of the same minors, with the signs of
the permutations that put the occupied slots in copy order. The sign
within a batch is already present in $U$. Within each row we may
reorder the occupied bits by batch by a signed basis permutation
local to that row. What remains is to interchange blocks belonging
to different rows and different batches. If $p_{ia}$ is the parity
of the occupied count in row $i$ of batch $a$, the remaining sign is
\[
       (-1)^{\sum_{i<j}\sum_{a>c}p_{ia}p_{jc}} .
\]
An earlier row in a later batch crosses a later row in an earlier
batch, with one inversion per occupied pair.

Supply the sign tensor on $d$ parity words of length $\ell$ at the
tip and send one word to each row, at cost at most $d\ell\log 2$.
Tensor this supply with $U^{\otimes\ell}$. A local map at each row
keeps the parity word matching its letters, thereby multiplying each
coefficient by the required sign. In characteristic two the sign is
one. In a fixed exterior sector the count in each slot fixes the total
weight on its row. We can therefore start with the weighted sector
of $U^{\otimes\ell}$, apply the sign correction, and project further
onto the prescribed slot counts. Maximizing over exterior sectors and
dividing by $\ell$ gives, in the limit,
\[
 b\left(\lambda(C_V)-d\sum_{j=1}^{q}\max(\theta_j,0)\right)
 \le \mathcal B_x(U)+d\log 2 .
\]
Here the exterior price is the limit along copy counts $b\ell$. The
subtraction accounts for the shift in Equation~\eqref{pack:letters}.
By Theorem~\ref{ext:star} the left side is at least
$d\mathcal F(1)-O(bM(1+d/q)\log n)$. Since $w\ge dm$ gives
$d/q\le m$, the error is $O(m^6)=o(d)$. Together with the cost of the
sign correction this yields $\mathcal B_x(U)\ge d\mathcal F(1)-O(d)$.

\smallskip
\noindent\emph{Compressing the row alphabet.}
The full alphabet uses $bq=\Theta(wm)$ bits per row. We retain the
letters $V=\{v:t(v)\ge-Aw\}$ for a sufficiently large absolute
integer $A$. Each retained letter contributes at least $e^{-Aw}$ to
the sum in Equation~\eqref{pack:free} at $r=1$, so
$|V|\le\exp(Aw+\mathcal F(1))$. Restricting the sum to $V$ also gives
Equation~\eqref{pack:partition}.

Project $U$ onto $V$ in every row to form $\widehat U$. In the full
format write $U$ as this projection plus $d$ tensors, one for each
possible first excluded row. On such a row one has
\[
 \sum_{t(v)<-Aw} e^{t(v)}
  \le e^{-Aw/2}\sum_v e^{t(v)/2}
  =\exp\big(-Aw/2+\mathcal F(2)/2\big).
\]
Choose $A$ so that the logarithm of the right side is at most
$\mathcal F(1)-w$. By transport from the tip, any tensor with an
excluded row has weighted price at most $d\mathcal F(1)-w$. Summing
the $d+1$ tensors costs at most $D_x\log(d+1)=O(m^{3/2})$ by
Equation~\eqref{pack:sum} and Lemma~\ref{price:diagonal}. Since
$w\ge dm$, the lower bound on $\mathcal B_x(U)$ forces
$\mathcal B_x(\widehat U)\ge d\mathcal F(1)-O(d)$. Removing the zero
coordinates outside $V$ preserves the sector prices by mutual
restriction. Transport from the tip gives the upper bound in
Equation~\eqref{pack:high}.

\smallskip
\noindent\emph{One binary tensor for all Vandermonde matrices.}
Replace the matrix parameters by indeterminates $z_1,\ldots,z_n$ in
slot order. For a tuple of retained row letters, let $S_e$ be the
occupied slots in copy $e$. The symbolic entry is zero if some
$|S_e|\ne n/2$. Otherwise it is the integral polynomial
\[
 P(z)=\pm\prod_{e=1}^b\
      \prod_{\substack{i,j\in S_e\\i<j}}(z_j-z_i).
\]
The sign reorders the occupied bits from slot order, with copy index
innermost, to copy order, with slots increasing. The total degree is
at most $D=bn^2$.

To encode the polynomial, attach an address $a_j\in\{0,\ldots,D\}$ to
slot $j$, carried on the row containing that slot. Attach one more
index $0\le s\le D^2$ to the last row. At these labels set the entry
of an integer tensor $T_0$ to
\begin{equation}\label{pack:univariate}
       [u^s]P(u^{a_1},\ldots,u^{a_n}).
\end{equation}
We claim every specialization $z_j=\xi_j$ is a rowwise restriction
of $T_0$ over $K$. Pick $u_0\in K^\times$ such that
$1,u_0,\ldots,u_0^D$ are distinct. Finitely many roots of unity are
forbidden, so this is possible in any infinite field. Let $L_a$ be
the Lagrange basis polynomial for node $u_0^a$ on this list. Send a
row basis vector to its retained letter, multiplied by
$\prod_{j\text{ in row}} L_{a_j}(\xi_j)$ and, on the last row, by
$u_0^s$. Summing over $s$ evaluates the polynomial in
Equation~\eqref{pack:univariate}, of degree at most $D^2$, at $u_0$.
Summing over the addresses interpolates separately in each $z_j$ and
recovers $P(\xi_1,\ldots,\xi_n)$. These maps work with the same
integer coefficients reduced in every characteristic.

The logarithm of the number of labels on any row is at most
$O(w)+q\log(D+1)+\log(D^2+1)=O(w)$, since $\log D=O(m)$. To compute
an entry from its row indices, first recover the retained letters.
Multiplying the penalty $-t(v)$ by $q$ gives nonnegative integer item
costs for the $bq$ bits, with total budget $Aqw$. A dynamic program
counts suffix completions at each remaining budget. These counts,
using at most $bq+1$ bits, allow lexicographic ranking and unranking
in time polynomial in $w$. For the recovered letters and addresses,
multiply at most $D$ binomials in Equation~\eqref{pack:univariate},
keeping degrees through $D^2$. Each intermediate coefficient has
absolute value at most $2^D$, and the exterior sign is computed by
counting inversions. All of these operations take polynomial time in
$w$.

Finally, append a sign $\varepsilon\in\{-1,1\}$ and a binary position
$0\le j\le D$ to the last leg. Set an entry to one precisely when the
corresponding coefficient of $T_0$ has sign $\varepsilon$ and its
absolute value has bit $j$ set. The last-leg map with multipliers
$\varepsilon 2^j$ recovers $T_0$ over every field. The enlarged
tensor has polynomial-time decidable zero-one entries and uses
$O(w)$ bits per row label. Pad each leg with zero coordinates to
$q_0=2^{\lceil C_0w\rceil}$ for an absolute integer $C_0$ large enough
to accommodate the labels and ensure $\log q_0\ge w$.
\end{proof}
